\documentclass[11pt]{amsart}

\usepackage[T1]{fontenc}
\usepackage{lmodern}
\usepackage{microtype}
\usepackage{mathtools,amssymb,amsthm}
\usepackage[hidelinks]{hyperref}

\hypersetup{
  pdftitle={A 25-vertex weak tensor product of alpha-free trees with an alpha-valuation}
}

\newtheorem{theorem}{Theorem}
\newtheorem{lemma}[theorem]{Lemma}
\newtheorem{proposition}[theorem]{Proposition}
\newtheorem{corollary}[theorem]{Corollary}
\theoremstyle{definition}
\newtheorem{fact}[theorem]{Fact}
\theoremstyle{remark}
\newtheorem*{remark}{Remark}

\newcommand{\wtp}{\mathbin{\overline{\otimes}}}
\newcommand{\Vsm}{V_{\mathrm{small}}}
\newcommand{\Vlg}{V_{\mathrm{large}}}

\title[A weak tensor product of $\alpha$-free trees with an $\alpha$-valuation]
{A 25-Vertex Weak Tensor Product of $\alpha$-Free Trees With an $\alpha$-Valuation}

\date{}

\subjclass[2020]{05C78, 05C05, 05B10}
\keywords{graceful labelling, $\alpha$-valuation, near $\alpha$-valuation,
gracious labelling, weak tensor product, counterexample}

\begin{document}

\begin{abstract}
Anglin, Huczynska and McCourt ask whether the weak tensor product of two graphs
that each admit a near $\alpha$-valuation but no $\alpha$-valuation must itself
admit no $\alpha$-valuation. The answer is no. Let $S_{3,2}$ be Rosa's
$7$-vertex spider, a tree of smallest size with no $\alpha$-valuation, carrying the
near $\alpha$-valuation printed in their paper. Its weak tensor product with
itself, taken with respect to the induced bipartitions, is a $25$-vertex
$36$-edge graph which admits an $\alpha$-valuation with boundary $\lambda=27$;
the valuation is written out in full, and the check is a pencil computation of
$36$ differences. The same conclusion holds for the non-isomorphic pair
$S_{3,2},S_{4,2}$, whose product on $32$ vertices and $48$ edges is exhibited
with boundary $\lambda=36$, so the question is not rescued by reading ``two
graphs'' strictly. Nothing published is contradicted. For \emph{connected}
products we record why a positive answer was plausible: the specific product
labelling of El-Zanati, Kenig and Vanden Eynden, on which the question rests,
is graceful but has no boundary as soon as neither factor labelling is an
$\alpha$-valuation, so it cannot itself witness such a counterexample.
\end{abstract}

\maketitle

\section{Introduction}

All graphs are finite and simple. Let $G$ have $q$ edges. Following Rosa
\cite{Rosa}, a \emph{$\beta$-valuation} (equivalently, a \emph{graceful
labelling}) of $G$ is an injection $b\colon V(G)\to\{0,1,\dots,q\}$ whose
multiset of edge differences $\{\,|b(u)-b(v)| : \{u,v\}\in E(G)\,\}$ is exactly
$\{1,2,\dots,q\}$. Quoting \cite{AHM} verbatim, its closing clause set inline
rather than displayed:

\begin{quote}
``An $\alpha$-valuation $b$ of a graph $G$ with $n$ edges is a
$\beta$-valuation which satisfies the following extra condition: there exists
some $x$ with $0\le x\le n$ such that for each edge $\{u,v\}\in E(G)$:
$b(v)\le x<b(u)$ or $b(u)\le x<b(v)$.''
\end{quote}

We call such an $x$ a \emph{boundary} and write $\lambda$ for it. Again
verbatim from \cite{AHM}, its two-item list set inline:

\begin{quote}
``A near $\alpha$-valuation $b$ of a graph $G$ with $n$ edges is a
$\beta$-valuation which satisfies the following extra condition: there is a
partition of $V(G)$ into $\Vsm$ and $\Vlg$ such that for each $v\in\Vsm$,
$b(v)<b(u)$ for all $u\in N(v)$, and for each $v\in\Vlg$, $b(v)>b(u)$ for all
$u\in N(v)$.''
\end{quote}

The notion has three published names: it is Cahit's 1980 notion
\cite{Cahit}, the \emph{gracious} labellings of Grannell, Griggs and Holroyd,
and the \emph{near $\alpha$-labelings} of El-Zanati, Kenig and Vanden Eynden
\cite{EKV}; \cite{AHM} records this coincidence: ``This concept was introduced
independently in [Cahit], [Grannell] and [El-Zanati], under different names.''

Finally, Snevily's \emph{weak tensor product} \cite{Snevily}, quoted from
\cite{AHM} with its display-math clause:

\begin{quote}
``The weak tensor product of two bipartite graphs $G$ and $H$ with partitions
$V_1$ and $V_2$ and $W_1$ and $W_2$ respectively is denoted $G\wtp H$. It is a
graph on the vertices $(V_1\times W_1)\cup(V_2\times W_2)$ where vertices
$(x,y)$ and $(u,v)$ are adjacent if and only if
$\{x,u\}\in E(G)$ and $\{y,v\}\in E(H)$.''
\end{quote}

Which bipartition is used is not free. El-Zanati, Kenig and Vanden Eynden
\cite[Thm.~7]{EKV}, restated in \cite{AHM}, states that if $G$ and $H$ admit
near $\alpha$-valuations $\gamma$ and $\delta$, then $G\wtp H$ \emph{taken with
respect to the vertex bipartitions induced by $\gamma$ and $\delta$} admits a
near $\alpha$-valuation, and \cite{AHM} adds: ``Note that it is essential to use
this specific choice of bipartitions.'' Throughout, $G\wtp H$ means the product
on
$(V(G)_{\mathrm{small}}\times V(H)_{\mathrm{small}})\cup
(V(G)_{\mathrm{large}}\times V(H)_{\mathrm{large}})$.

\subsection*{The question}
The second of the four unnumbered items in the ``Further Work'' section of
\cite{AHM} reads verbatim and in full, with its cross-reference rendered as the
ordinary citation \cite[Thm.~7]{EKV}:

\begin{quote}
``From Theorem \cite[Thm.~7]{EKV}, the weak tensor product of any two
graphs with near-$\alpha$ valuations is a graph with a near $\alpha$-valuation.
For two graphs which each possess a near $\alpha$ valuation but no
$\alpha$-valuation, will their weak tensor product have no $\alpha$-valuation?''
\end{quote}

It is a question, not a numbered conjecture: it carries no theorem environment
and no number. Read as the universally quantified statement it invites:

\begin{quote}
$(\ast)$ \emph{For all bipartite $G,H$ such that $G$ admits a near
$\alpha$-valuation and no $\alpha$-valuation, and likewise $H$, the product
$G\wtp H$ taken with respect to the bipartitions induced by those valuations
admits no $\alpha$-valuation.}
\end{quote}

\begin{theorem}\label{thm:main}
$(\ast)$ is false, and the question quoted above is answered in the
negative. Let $S_{3,2}$ be the tree obtained from $K_{1,3}$ by subdividing each
edge once, with the near $\alpha$-valuation $\gamma$ printed in
\textup{\cite{AHM}}. Then $S_{3,2}$ admits no $\alpha$-valuation, while
$K=S_{3,2}\wtp S_{3,2}$ --- a connected bipartite graph on $25$ vertices and
$36$ edges with parts of sizes $16$ and $9$ --- admits the $\alpha$-valuation
displayed in \S\ref{sec:witness}, with boundary $\lambda=27$.
\end{theorem}

The same holds for a non-isomorphic pair (\S\ref{sec:distinct}), so the
question is not rescued by insisting that ``two graphs'' means two distinct
graphs. \S\ref{sec:why} records one reason a positive answer was plausible: when
the product is connected, the particular product labelling of \cite{EKV}, on
which the question rests, cannot itself produce a counterexample.

Theorem~7 of \cite{EKV} (near $\alpha$ in, near $\alpha$ out) and Snevily's
theorem \cite{Snevily} ($\alpha$ in, $\alpha$ out) both remain true; what fails
is only the hoped-for propagation of $\alpha$-\emph{freeness}. See
\S\ref{sec:scope} for the limits of what is claimed.

\section{The factor, and its $\alpha$-freeness}\label{sec:factor}

Write $S_{3,2}$ for the spider with centre $c$, subdivision vertices
$m_1,m_2,m_3$ and leaves $l_1,l_2,l_3$, and edges
\[
 E(S_{3,2})=\{\,cm_1,\;cm_2,\;cm_3,\;l_1m_1,\;l_2m_2,\;l_3m_3\,\},
\]
so $|V|=7$ and $q=6$. Its unique bipartition is
$\{c,l_1,l_2,l_3\}\mid\{m_1,m_2,m_3\}$. \cite{AHM} states that ``a tree of
smallest size which does not possess an $\alpha$-valuation is given in
\cite{Rosa}; it is the 7-vertex graph $S_{3,2}$'', and prints a near
$\alpha$-valuation of it. Read off that figure, the valuation is
\begin{equation}\label{eq:gamma}
 \gamma(c)=0,\qquad
 (\gamma(m_i),\gamma(l_i))_{i=1,2,3}=(5,4),\,(3,1),\,(6,2).
\end{equation}
Its six edge differences are $5,3,6$ (centre edges) and $1,2,4$ (leaf edges),
i.e.\ exactly $\{1,\dots,6\}$, so $\gamma$ is a $\beta$-valuation; and with
\[
 \Vsm=\{c,l_1,l_2,l_3\}\ \ (\text{labels }0,4,1,2),
 \quad
 \Vlg=\{m_1,m_2,m_3\}\ (\text{labels }5,3,6),
\]
every vertex of $\Vsm$ is below and every vertex of $\Vlg$ above all its
neighbours, so $\gamma$ is a near $\alpha$-valuation. It is not an
$\alpha$-valuation: the two label sets interleave, since $4\in\gamma(\Vsm)$
exceeds $3\in\gamma(\Vlg)$.

The following is standard; both later arguments turn on it.

\begin{fact}\label{fact:0}
Let $T$ be a connected bipartite graph and let $b$ be a near
$\alpha$-valuation of $T$ with classes $\Vsm,\Vlg$.
\begin{enumerate}
\item[(i)] $b$ is an $\alpha$-valuation if and only if
$\max b(\Vsm)<\min b(\Vlg)$, in which case $\lambda=\max b(\Vsm)$ is a
boundary.
\item[(ii)] If in addition $|V(T)|=q+1$, where $q=|E(T)|$, then every
$\alpha$-valuation $b'$ of $T$ with boundary $x$ is a bijection onto
$\{0,\dots,q\}$, the two sides of the unique bipartition of $T$ are
$\{v:b'(v)\le x\}$ and $\{v:b'(v)>x\}$, and the first receives exactly the
labels $\{0,1,\dots,x\}$.
\end{enumerate}
\end{fact}

\begin{proof}
(i) If $\max b(\Vsm)<\min b(\Vlg)$ then $x=\max b(\Vsm)$ straddles every edge,
because every edge of $T$ joins $\Vsm$ to $\Vlg$. Conversely suppose $b$ has
boundary $x$. The sets $L=\{v:b(v)\le x\}$ and $U=\{v:b(v)>x\}$ are
independent, since every edge straddles $x$, so $\{L,U\}$ is a bipartition of
$T$; a connected bipartite graph has only one, so $\{L,U\}=\{\Vsm,\Vlg\}$. As
every edge runs from $\Vsm$ up to $\Vlg$ we cannot have $L=\Vlg$, so
$L=\Vsm$ and $\max b(\Vsm)\le x<\min b(\Vlg)$.

(ii) Let $b'$ be any $\alpha$-valuation of $T$ with boundary $x$. It injects
into $\{0,\dots,q\}$ and $|V(T)|=q+1$, so it is a bijection. Put
$L=\{v:b'(v)\le x\}$ and $U=\{v:b'(v)>x\}$. Every edge straddles $x$, so both
are independent and $\{L,U\}$ is a bipartition of $T$, hence \emph{the}
bipartition. Finally $b'(L)\subseteq\{0,\dots,x\}$ and $|L|=x+1$ because $b'$
is a bijection, so $b'(L)=\{0,\dots,x\}$.
\end{proof}

\begin{proposition}\label{prop:s32}
$S_{3,2}$ has no $\alpha$-valuation.
\end{proposition}

\begin{proof}
Suppose $b$ is one, with boundary $x$. By Fact~\ref{fact:0}(ii), $b$ is a bijection
onto $\{0,\dots,6\}$, the low side is one of the two bipartition classes, and
it receives an initial segment of labels. Every edge runs low to high, so
\[
 \sum_{\{u,v\}\in E}\bigl|b(u)-b(v)\bigr|
 \;=\;\sum_{v\in\text{high}}\deg(v)\,b(v)-\sum_{v\in\text{low}}\deg(v)\,b(v)
 \;=\;1+\dots+6=21,
\]
with $\deg c=3$, $\deg m_i=2$, $\deg l_i=1$.

\emph{Case A: low $=\{c,l_1,l_2,l_3\}$}, receiving $\{0,1,2,3\}$, so the
$m_i$ receive $\{4,5,6\}$. The sum is
$2(4{+}5{+}6)-\bigl(3b(c)+(6-b(c))\bigr)=24-2b(c)$, which is even.

\emph{Case B: low $=\{m_1,m_2,m_3\}$}, receiving $\{0,1,2\}$, so $c$ and the
leaves receive $\{3,4,5,6\}$. The sum is
$\bigl(3b(c)+(18-b(c))\bigr)-2(0{+}1{+}2)=2b(c)+12$, again even.

In both cases the sum is even, but $21$ is odd.
\end{proof}

No computation is used; in particular the appeal to \cite{Rosa} is not needed.

\section{The product, and the $\alpha$-valuation}\label{sec:witness}

\subsection*{The graph $K$, described without reference to the product}
$K$ can be built directly. Take a hub $h$; a $3\times3$ grid of vertices
$Q_{ij}$ ($1\le i,j\le3$); three ``row'' vertices $R_1,R_2,R_3$; three
``column'' vertices $C_1,C_2,C_3$; and nine pendants $P_{ij}$. Join $h$ to all
nine $Q_{ij}$; join $R_i$ to $Q_{i1},Q_{i2},Q_{i3}$; join $C_j$ to
$Q_{1j},Q_{2j},Q_{3j}$; and join $P_{ij}$ to $Q_{ij}$. Then
$|V(K)|=1+9+3+3+9=25$ and $|E(K)|=9+9+9+9=36$. The degree sequence is
$9^1\,4^9\,3^6\,1^9$, $K$ is connected, and
\[
 P=\{h\}\cup\{R_i\}\cup\{C_j\}\cup\{P_{ij}\}\ (|P|=16),
 \qquad
 Q=\{Q_{ij}\}\ (|Q|=9)
\]
is its unique bipartition.

This is $S_{3,2}\wtp S_{3,2}$ with $\gamma=\delta$ as in \eqref{eq:gamma}.
Indeed, by the definition quoted in \S1 the product lives on
$\bigl(\{c,l_1,l_2,l_3\}^2\bigr)\cup\bigl(\{m_1,m_2,m_3\}^2\bigr)$, i.e.\ on
$16+9=25$ vertices, and $(x,y)\sim(u,v)$ iff $\{x,u\}$ and $\{y,v\}$ are edges
of $S_{3,2}$. Writing $xy$ for $(x,y)$:

\begin{center}
\begin{tabular}{llll}
\hline
vertex of $K$ & neighbours & degree & r\^ole\\
\hline
$cc$        & all nine $m_im_j$ & $9$ & $h$\\
$l_ic$      & $m_im_1,m_im_2,m_im_3$ & $3$ & $R_i$\\
$cl_j$      & $m_1m_j,m_2m_j,m_3m_j$ & $3$ & $C_j$\\
$l_il_j$    & $m_im_j$ only & $1$ & $P_{ij}$\\
$m_im_j$    & $cc,\;cl_j,\;l_ic,\;l_il_j$ & $4$ & $Q_{ij}$\\
\hline
\end{tabular}
\end{center}

\subsection*{The valuation}
Set $\lambda=27$ and define $b$ on $V(K)$ by the two tables

\begin{align*}
 P\text{-side } (=\{c,l_1,l_2,l_3\}^2,\ 16 \text{ labels}):\quad
 & cl_2=0,\ l_1l_2=1,\ l_1c=2,\ cl_1=3,\\
 & l_2l_3=4,\ l_2l_2=5,\ l_3l_1=7,\ l_3c=8,\\
 & l_1l_1=9,\ cl_3=15,\ l_2c=17,\ l_3l_3=18,\\
 & l_3l_2=20,\ l_2l_1=21,\ l_1l_3=23,\ cc=\mathbf{27};\\[2pt]
 Q\text{-side } (=\{m_1,m_2,m_3\}^2,\ 9 \text{ labels}):\quad
 & m_1m_1=\mathbf{28},\ m_3m_3=29,\ m_3m_2=30,\\
 & m_3m_1=31,\ m_2m_2=32,\ m_2m_3=33,\\
 & m_2m_1=34,\ m_1m_3=35,\\
 & m_1m_2=36.
\end{align*}

\begin{proof}[Proof of Theorem~\ref{thm:main}]
Proposition~\ref{prop:s32} and \S\ref{sec:factor} give the hypotheses: $\gamma$
is a near $\alpha$-valuation of $S_{3,2}$ and $S_{3,2}$ has no
$\alpha$-valuation. It remains to check that $b$ above is an
$\alpha$-valuation of $K$.

\emph{Injectivity and range.} The $25$ values
\[
 0,1,2,3,4,5,7,8,9,15,17,18,20,21,23,27\ \bigm|\ 28,29,\dots,36
\]
are pairwise distinct and lie in $\{0,\dots,36\}=\{0,\dots,|E(K)|\}$. Exactly
$37-25=12$ labels are unused, namely
$6,10,11,12,13,14$, $16,19,22,24,25,26$.

\emph{The boundary.} $\max_P b=27$ (attained at the hub $cc$) and
$\min_Q b=28$, while every edge of $K$ joins $P$ to $Q$. Hence
$b(u)\le 27<b(v)$ on all $36$ edges, so $\lambda=27$ is a boundary.

\emph{Gracefulness.} In the r\^oles of the intrinsic description,
$Q_{11}=28$, $Q_{12}=36$, $Q_{13}=35$, $Q_{21}=34$, $Q_{22}=32$, $Q_{23}=33$,
$Q_{31}=31$, $Q_{32}=30$, $Q_{33}=29$, and the $36$ differences split into
four blocks:
\begin{align*}
 h=cc=27     &\to Q_{11},\dots,Q_{33}: & &1,\,9,\,8,\,7,\,5,\,6,\,4,\,3,\,2\\
 R_1=l_1c=2  &\to Q_{11},Q_{12},Q_{13}: & &26,\,34,\,33\\
 R_2=l_2c=17 &\to Q_{21},Q_{22},Q_{23}: & &17,\,15,\,16\\
 R_3=l_3c=8  &\to Q_{31},Q_{32},Q_{33}: & &23,\,22,\,21\\
 C_1=cl_1=3  &\to Q_{11},Q_{21},Q_{31}: & &25,\,31,\,28\\
 C_2=cl_2=0  &\to Q_{12},Q_{22},Q_{32}: & &36,\,32,\,30\\
 C_3=cl_3=15 &\to Q_{13},Q_{23},Q_{33}: & &20,\,18,\,14\\
 P_{ij}      &\to Q_{ij}: & &19,\,35,\,12,\ \ 13,\,27,\,29,\ \ 24,\,10,\,11
\end{align*}
(the last line lists $P_{11},P_{12},P_{13}$, then $P_{21},P_{22},P_{23}$, then
$P_{31},P_{32},P_{33}$, whose labels are $9,1,23$, then $21,5,4$, then
$7,20,18$). Sorting the union and tagging each value with the block it came
from:
\[
\begin{array}{l}
 1_h\ \ 2_h\ \ 3_h\ \ 4_h\ \ 5_h\ \ 6_h\ \ 7_h\ \ 8_h\ \ 9_h\ \ 10_P\ \ 11_P\
 \ 12_P\\
 13_P\ \ 14_C\ \ 15_R\ \ 16_R\ \ 17_R\ \ 18_C\ \ 19_P\ \ 20_C\ \ 21_R\ \ 22_R\
 \ 23_R\ \ 24_P\\
 25_C\ \ 26_R\ \ 27_P\ \ 28_C\ \ 29_P\ \ 30_C\ \ 31_C\ \ 32_C\ \ 33_R\ \ 34_R\
 \ 35_P\ \ 36_C
\end{array}
\]
which is $1,2,\dots,36$, each exactly once. So $b$ is a $\beta$-valuation, and
with the boundary above it is an $\alpha$-valuation.
\end{proof}

\begin{remark}
Two independent consistency checks on the printed numbers. First, summing
$b(\text{high end})-b(\text{low end})$ over all edges gives
$\sum_{v\in Q}\deg(v)b(v)-\sum_{v\in P}\deg(v)b(v)=666=1+2+\dots+36$. Second,
the edge carrying difference $36$ is $cl_2=0\sim m_1m_2=36$, and those two
vertices \emph{are} adjacent, since $\{c,m_1\}$ and $\{l_2,m_2\}$ are edges of
$S_{3,2}$.
\end{remark}

\begin{corollary}
$\alpha$-freeness is not preserved by the weak tensor product, even for trees
as factors and even when the product is taken with respect to the bipartitions
induced by near $\alpha$-valuations of the factors.
\end{corollary}

\section{The product labelling $\sigma$}\label{sec:why}

When the product is connected, the labelling of $G\wtp H$ on which the question
rests cannot itself witness a counterexample. Let
$A=V(G)_{\mathrm{small}}$, $B=V(G)_{\mathrm{large}}$, $m=|E(G)|$,
$g_A=\max\gamma(A)$, $g_B=\min\gamma(B)$, and let $C,D,n,h_C,h_D$ be the
corresponding data for $H$ and $\delta$. The labelling $\sigma$ of $G\wtp H$
used by \cite{EKV} and reproduced in \cite{AHM} is
\begin{align*}
 \sigma(a,c)&=m\,\delta(c)+\gamma(a) & &(a\in A,\ c\in C),\\
 \sigma(b,d)&=m\,(\delta(d)-1)+\gamma(b) & &(b\in B,\ d\in D),
\end{align*}
and is a $\beta$-valuation of $G\wtp H$, with edge differences $m(f-1)+e$ for
$e\in\{1,\dots,m\}$, $f\in\{1,\dots,n\}$.

\begin{lemma}\label{lem:sigma}
Suppose $G\wtp H$ is connected. Then $\sigma$ is an $\alpha$-valuation of
$G\wtp H$ if and only if
\[
 m\,(h_C-h_D+1)\;<\;g_B-g_A .
\]
\end{lemma}

\begin{proof}
Every edge of $G\wtp H$ runs from $A\times C$ to $B\times D$, and $\sigma$ is
a near $\alpha$-valuation with those classes \cite[Thm.~7]{EKV}. By
Fact~\ref{fact:0}(i), $\sigma$ is an $\alpha$-valuation
iff $\max\sigma(A\times C)<\min\sigma(B\times D)$. Now
$\max\sigma(A\times C)=m\,h_C+g_A$ and
$\min\sigma(B\times D)=m\,(h_D-1)+g_B$, and rearranging gives the claim.
\end{proof}

\begin{corollary}\label{cor:obstruction}
Suppose $G\wtp H$ is connected. If neither $\gamma$ nor $\delta$ is an
$\alpha$-valuation, then $\sigma$ is not an $\alpha$-valuation of $G\wtp H$ ---
for every such choice of $\gamma$ and $\delta$ and in either coordinate order.
Conversely, if both $\gamma$ and $\delta$ are $\alpha$-valuations then $\sigma$
is one, which is Snevily's theorem \textup{\cite{Snevily}} for $\wtp$.
\end{corollary}

\begin{proof}
Labels are distinct, so by Fact~\ref{fact:0}(i) exactly one of $g_A<g_B$,
$g_A>g_B$ holds, and $\gamma$ is $\alpha$ iff the first does; likewise for
$\delta$. If neither is $\alpha$ then $g_B-g_A\le-1$ and $h_C-h_D\ge1$, so
$m(h_C-h_D+1)\ge 2m\ge2>-1\ge g_B-g_A$ and Lemma~\ref{lem:sigma} fails. If
both are $\alpha$ then $g_B-g_A\ge1$ and $h_C-h_D\le-1$, so
$m(h_C-h_D+1)\le0<1\le g_B-g_A$.
\end{proof}

Connectedness of $G\wtp H$ is used in both statements, through the uniqueness of
the bipartition in Fact~\ref{fact:0}(i); we do not settle what $\sigma$ does when
$G\wtp H$ is disconnected. Both products exhibited in this note are connected.

For our $K$, $m=6$, $h_C=g_A=4$ and $h_D=g_B=3$, so the criterion reads
$12<-1$: $\sigma$ gives labels with $\max_P\sigma=28>15=\min_Q\sigma$, graceful
but with no boundary. The witness $b$ of \S\ref{sec:witness} is therefore not
$\sigma$; it is also not \emph{separable}, i.e.\ not of the form
$\varphi(x,y)=p(x)+q(y)$, since $b(cl_1)-b(cl_2)=3$ while
$b(l_1l_1)-b(l_1l_2)=8$, which separability would force to be equal. We do not
claim that a counterexample must be non-separable.

\section{Two distinct graphs}\label{sec:distinct}

Theorem~7 of \cite{EKV}, on which the bullet explicitly rests, has no
distinctness or non-isomorphism hypothesis, so Theorem~\ref{thm:main} already
answers it. The strict reading of ``two graphs'' does not rescue it either.

Let $S_{4,2}$ be the comet with centre $c$, subdivision vertices
$a_0,\dots,a_3$ and leaves $b_0,\dots,b_3$, edges $ca_j$ and $b_ja_j$; so
$|V|=9$, $q=8$, and the unique bipartition is
$\{c,b_0,\dots,b_3\}\mid\{a_0,\dots,a_3\}$. The labelling
\begin{equation}\label{eq:delta}
 \delta(c)=0,\qquad
 (\delta(a_j),\delta(b_j))_{j=0,1,2,3}=(8,4),\,(7,2),\,(6,5),\,(3,1)
\end{equation}
is a bijection onto $\{0,\dots,8\}$ with edge differences $8,7,6,3$ from the
centre edges and $4,5,1,2$ from the leaf edges, i.e.\ exactly $\{1,\dots,8\}$.
It is a near $\alpha$-valuation, with $\Vsm$ the set of $c$ and the four
leaves; and it is not an $\alpha$-valuation, since $5=\delta(b_2)$ exceeds
$3=\delta(a_3)$.

\begin{proposition}\label{prop:s42}
$S_{4,2}$ has no $\alpha$-valuation.
\end{proposition}

\begin{proof}
As in Proposition~\ref{prop:s32}, an $\alpha$-valuation $b$ is a bijection onto
$\{0,\dots,8\}$, the low side is a bipartition class receiving an initial
segment, and the edge differences sum to $1+\dots+8=36$. Here parity alone does
not decide, so we push one step further.

\emph{Case A: low $=\{c,b_0,\dots,b_3\}$}, receiving $\{0,1,2,3,4\}$; the $a_j$
receive $\{5,6,7,8\}$. The sum is
$2(5{+}6{+}7{+}8)-\bigl(4b(c)+(10-b(c))\bigr)=42-3b(c)$, so $b(c)=2$ is
forced. Then the centre edges contribute the differences
$\{5{-}2,6{-}2,7{-}2,8{-}2\}=\{3,4,5,6\}$, so the four leaf edges must
contribute $\{1,2,7,8\}$. The leaves receive $\{0,1,3,4\}$ and each is matched
with a distinct $a_j$. The only pair in
$\{5,6,7,8\}\times\{0,1,3,4\}$ with difference $1$ is $(5,4)$, and the only
ones with difference $2$ are $(5,3)$ and $(6,4)$; using $(5,4)$ consumes both
$5$ and $4$, so differences $1$ and $2$ cannot both occur.

\emph{Case B: low $=\{a_0,\dots,a_3\}$}, receiving $\{0,1,2,3\}$; then $c$ and
the leaves receive $\{4,\dots,8\}$. The sum is
$\bigl(4b(c)+(30-b(c))\bigr)-2(0{+}1{+}2{+}3)=3b(c)+18$, so $b(c)=6$ is
forced. The centre edges contribute $\{6,5,4,3\}$, so again the leaf edges must
contribute $\{1,2,7,8\}$. The leaves receive $\{4,5,7,8\}$; the only pair with
difference $1$ is $(4,3)$ and the only ones with difference $2$ are $(4,2)$
and $(5,3)$, and $(4,3)$ consumes both $4$ and $3$. Same contradiction.
\end{proof}

\begin{theorem}\label{thm:distinct}
$(\ast)$ fails on a non-isomorphic pair. With $\gamma$ as in
\eqref{eq:gamma} on $G=S_{3,2}$ (relabel $m_i\mapsto a_{i-1}$,
$l_i\mapsto b_{i-1}$) and $\delta$ as in \eqref{eq:delta} on $H=S_{4,2}$, the
product $L=G\wtp H$ is a connected bipartite graph on $32$ vertices and $48$
edges with parts of sizes $20$ and $12$, and
\begin{align*}
 P\ (\text{low, }20):\quad & b_0b_1=0,\ b_0c=1,\ cb_3=2,\ cb_0=3,\ cb_2=4,\\
 & b_2b_0=5,\ b_2c=9,\ b_2b_2=10,\ b_1c=17,\ b_1b_3=18,\\
 & b_1b_1=19,\ b_1b_2=20,\ b_0b_0=22,\ b_1b_0=23,\ b_2b_1=24,\\
 & cb_1=25,\ b_2b_3=26,\ b_0b_2=27,\ b_0b_3=28,\ cc=\mathbf{36};\\[2pt]
 Q\ (\text{high, }12):\quad
 & a_0a_0=\mathbf{37},\ a_1a_1=38,\ a_2a_2=39,\ a_2a_3=40,\\
 & a_2a_1=41,\ a_2a_0=42,\ a_1a_3=43,\ a_1a_2=44,\\
 & a_1a_0=45,\ a_0a_3=46,\ a_0a_2=47,\ a_0a_1=48
\end{align*}
is an $\alpha$-valuation of $L$ with boundary $\lambda=36$.
\end{theorem}

\begin{proof}
Here $(x,y)\sim(u,v)$ iff $\{x,u\}\in E(G)$ and $\{y,v\}\in E(H)$, on
\[
 (\{c,b_0,b_1,b_2\}\times\{c,b_0,\dots,b_3\})\cup
 (\{a_0,a_1,a_2\}\times\{a_0,\dots,a_3\}),
\]
i.e.\ on $20+12=32$ vertices; the
degrees are $12$ at $cc$, $3$ at each $cb_j$, $4$ at each $b_ic$, $1$ at each
$b_ib_j$, so $|E(L)|=12+4\cdot3+3\cdot4+12=48$, and $L$ is connected. The $32$
labels above are distinct and lie in $\{0,\dots,48\}$; $\max_P=36<37=\min_Q$
and every edge runs $P\to Q$, so $\lambda=36$ straddles all $48$ edges. That
the $48$ differences are exactly $\{1,\dots,48\}$ is a finite check on the
printed labels; as a consistency test,
$\sum_{v\in Q}\deg(v)b(v)-\sum_{v\in P}\deg(v)b(v)=1176=1+2+\dots+48$.
Propositions~\ref{prop:s32} and~\ref{prop:s42} supply the hypotheses.
\end{proof}

\section{Relation to known results}

Every weak-tensor-product theorem recorded in Gallian's survey \cite{Gallian}
requires the factors to \emph{have} an ($\alpha$ or near-$\alpha$) labelling:
Snevily \cite{Snevily}, El-Zanati--Kenig--Vanden Eynden \cite[Thm.~7]{EKV}, and
the generalisation of L\'opez and Muntaner-Batle \cite{LM}. We found no
statement in either direction there about $\wtp$ applied to $\alpha$-free
factors, which is exactly the gap the question of \cite{AHM} occupies.

The nearest results should be recorded, since the phenomenon --- an
$\alpha$-free tree sitting inside an $\alpha$-labelled multi-copy host --- is
\emph{not} new. Ahmed and Snevily
\cite{AS} prove that for every comet $T$ either $T$ has an $\alpha$-labelling or
there is a graph $H_T$ with an $\alpha$-labelling that decomposes into two
edge-disjoint copies of $T$; the factors used here, $S_{3,2}$ and $S_{4,2}$, are
comets, so their result already applies to them. Their host is not a weak tensor
product, so it does not bear on $(\ast)$. El-Zanati, Fu and Shiue \cite{EFS}
prove Snevily's conjecture that every bipartite graph has a finite
$\alpha$-labelling number, i.e.\ that some number of edge-disjoint copies of it
tile a graph carrying an $\alpha$-labelling; again the host is not a weak tensor
product, and nothing there decides $(\ast)$. And L\'opez and Muntaner-Batle
\cite{LM}, which \cite{AHM} does not cite, enlarges the classes of graphs
admitting $\alpha$-, near-$\alpha$- and bigraceful labellings --- but its
$\alpha$-conclusion theorem requires $\alpha$-labellings on both factors, and its
near-$\alpha$ theorem concludes only near $\alpha$. Anglin, Huczynska and McCourt
note a further adjacency: ``However as shown in [Barrientos], graphs which
contain this as a subgraph may still have $\alpha$-valuations.''

What settles the question of \cite{AHM} is that the product construction the
question itself names is such a host. We do not claim that $\wtp$ never yields
$\alpha$-free graphs from $\alpha$-free factors: the two products above are
single instances.

\section{Scope}\label{sec:scope}

\begin{itemize}
\item What is refuted is $(\ast)$, the universally quantified reading of an
interrogative bullet. \cite{AHM} poses no conjecture there, and we
do not claim it does. Nothing published is contradicted; Theorem~7 of
\cite{EKV} and Snevily's theorem both stand, and Lemma~\ref{lem:sigma}
recovers the latter for connected products.
\item No general theorem in the opposite direction is proved: we do \emph{not}
claim that every such product has an $\alpha$-valuation.
\item Nothing is claimed about $\alpha$-free near-$\alpha$ trees beyond the two
factors used here, nor about non-tree $\alpha$-free bipartite factors, nor about
minimality: we do not claim $K$ is the smallest counterexample.
\item One residual is definitional rather than arithmetical: that ``near
$\alpha$-valuation'' and ``weak tensor product'' are read as \cite{AHM}
intends. Every quotation above is verbatim from the \LaTeX{} e-print of
\cite{AHM}, with no word added, dropped or altered; the only editing is of
formatting, in that the source's displayed formulae and its two-item list in the
near $\alpha$-valuation definition are set inline here, and its citation and
cross-reference macros are rendered as bracketed references. The reading was
pinned by reproducing that paper's own printed valuation \eqref{eq:gamma} and its
stated properties, but a definitional misreading is not something arithmetic can
exclude.
\item The accompanying program \texttt{verify.py}, with its recorded output
\texttt{verify.output.txt}, re-checks the objects printed above and the finite
claims made about them. It does not search for further counterexamples, does not
address minimality, and establishes nothing about factors on more than $9$
vertices.
\end{itemize}

\begin{thebibliography}{9}

\bibitem{AHM}
J.~Anglin, S.~Huczynska, and M.~McCourt,
\emph{Graph labellings and external difference families},
arXiv:2603.05662v1, 2026.
\href{https://arxiv.org/abs/2603.05662}{arXiv:2603.05662}.

\bibitem{AS}
T.~Ahmed and H.~Snevily,
\emph{The $\alpha$-labeling number of comets is 2},
Bull. Inst. Combin. Appl. \textbf{72} (2014), 25--40.

\bibitem{Cahit}
I.~Cahit,
\emph{Another formulation of the graceful tree conjecture},
Research Report CORR 80-47, Department of Combinatorics and Optimization,
University of Waterloo, 1980, 1--6.

\bibitem{EFS}
S.~I. El-Zanati, H.-L. Fu, and C.-L. Shiue,
\emph{On the $\alpha$-labeling number of bipartite graphs},
Discrete Math. \textbf{214} (2000), 241--243.
\href{https://doi.org/10.1016/S0012-365X(99)00313-1}{doi:10.1016/S0012-365X(99)00313-1}.

\bibitem{EKV}
S.~I. El-Zanati, M.~J. Kenig, and C.~Vanden Eynden,
\emph{Near $\alpha$-labelings of bipartite graphs},
Australas. J. Combin. \textbf{21} (2000), 275--285.

\bibitem{Gallian}
J.~A. Gallian,
\emph{A dynamic survey of graph labeling},
Electron. J. Combin., Dynamic Survey DS6; edition of 27 January 2026,
\S3.1.

\bibitem{LM}
S.~C. L\'opez and F.~A. Muntaner-Batle,
\emph{A new application of the $\otimes_h$-product to $\alpha$-labelings},
Discrete Math. \textbf{338} (2015), no.~6, 839--843.
\href{https://doi.org/10.1016/j.disc.2015.01.007}{doi:10.1016/j.disc.2015.01.007};
\href{https://arxiv.org/abs/1407.8537}{arXiv:1407.8537}.

\bibitem{Rosa}
A.~Rosa,
\emph{On certain valuations of the vertices of a graph},
Theory of Graphs (Internat.\ Sympos., Rome, 1966), Gordon and Breach, 1967,
349--355.

\bibitem{Snevily}
H.~S. Snevily,
\emph{New families of graphs that have $\alpha$-labelings},
Discrete Math. \textbf{170} (1997), 185--194.

\end{thebibliography}

\end{document}
